\documentclass[10pt,a4paper]{amsart}
\usepackage[margin=19mm]{geometry}
\usepackage{amsmath,amssymb,mathtools}
\usepackage[scr=rsfs,cal=euler]{mathalfa}
\usepackage{xfrac}
\usepackage[unicode,hidelinks,bookmarks=false]{hyperref}
\newcommand{\ie}{i.e.\ }
\newcommand{\cf}{cf.\ }
\newcommand{\resp}{respectively\ }
\newcommand{\Subsection}{\subsection}
\providecommand{\nobreakdash}{\nobreak\hbox{-}}
\setlength{\emergencystretch}{2em}
\multlinegap0pt
\usepackage{amssymb}
\newtheorem{lm}{Lemma}
\newtheorem{prop}{Proposition}
\title{Hereditarily and Super Bassian Modules \\ over Certain Rings}
\author{Peter Vassilev Danchev, Truong Cong Quynh, Jan \v{Z}emli\v{c}ka}
\date{Source: arXiv:2603.24188v1 (CC BY 4.0)}
\begin{document}
\maketitle

\noindent\emph{Source and licence.} Hereditarily and Super Bassian Modules \\ over Certain Rings. Version arXiv:2603.24188v1 (CC BY 4.0), distributed by arXiv under the Creative Commons Attribution 4.0 International licence, http://creativecommons.org/licenses/by/4.0/, as recorded in the arXiv licence metadata for this exact version and shown on the abstract page https://arxiv.org/abs/2603.24188v1. Full licence notice: \texttt{licenses/arxiv-2603.24188-CC-BY-4.0.txt}.\par\smallskip

\noindent\textit{Editorial scope.} Counted unit: the article's prop prop (source0.tex lines 180--186). The article's complete original proof of it is delivered with it (source0.tex lines 188--194, one \texttt{proof} environment of 7 lines). No mathematics is written, repaired or re-derived: the statement and the proof are the article's own lines, in the article's own notation, and no retained line is substituted or re-flowed.  Retained uncounted as the source-local context the proof needs: lines 164--169.  Nothing was omitted from any retained span, so the package carries no omission record.  No retained line contains a citation, so the wrapper carries no bibliography and no bibliography entry.  No retained line contains a figure environment, an image-inclusion command or drawing code, so no image is delivered and the package declares no asset dependency.  Editorial formatting only, in addition: an AMS \texttt{amsart} preamble that restates the article's own declarations and macros used by the retained lines, namely the environments \texttt{lm}, \texttt{prop} on separate counters, together with the standard packages those lines need.  The retained environments print in this document as Lm 1, Proposition 1 in order of appearance, whereas the article prints the same items with the numbering of its own full document.  The complete native source of the article as distributed in the arXiv e-print for this exact version is bundled unchanged as \texttt{sources/197128/source0.tex}.
\begin{lm}\label{perfect} Let $M$ be an infinitely generated module over a right perfect ring. Then,
\begin{enumerate}
\item there exists a simple module $S$ and a submodule $N$ such that $M/N\cong S^{(\omega)}$;
\item M is not super Bassian.
\end{enumerate}
\end{lm}

\begin{prop} Let $M$ be a right module over a right Artinian ring $R$. Then, the following three conditions are equivalent:
\begin{enumerate}
\item $M$ is super Bassian.
\item $M$ is finitely generated.
\item $M$ is a q.f.d. module.
\end{enumerate}	
\end{prop}

\begin{proof} (1) $\Rightarrow$ (2). It follows from Lemma~\ref{perfect}(2)

(2) $\Rightarrow$ (1) and (3). Since a finitely generated module over an Artinian ring is necessarily of finite length, it is a super Bassian q.f.d. module.
	
(3) $\Rightarrow$ (2). Assume that $M$ is an infinitely generated module. Then by
Lemma~\ref{perfect}(1), there exist a simple module \( S \) and a submodule \( N \) of $M$ such that \(M/N \cong S^{(\omega)}\), which is a contradiction to the finite Goldie dimension of $M/N$, as expected.
\end{proof}

\end{document}
